\chapter{線積分}
    \section{定義}
        $\realNumbers^n$上の有界関数$f$と区分的に滑らかな有向曲線$C$に対して$x_d\;(d=1,\dotsc,n)$軸方向の線積分を以下のようにして定義する。
        \par
        $C$上にその向きに沿って$n+1$個の分点\quad$\text{始点} = P_0, P_1, \dotsc, P_n = \text{終点}$\quad をとる。
        添字の意味で隣接する2点$P_{i-1},P_i$は重ならないようにする。さらに添字の順番が曲線の向きと逆行しないようにする。
        曲線が自己交差する点では複数の点が重なることがあるが、それは構わない。
        これらの点による$C$の分割を$\Delta$と表し、その「幅」を$|\Delta| \coloneqq \max_{i\in\range{1}{n}}\norm{P_{i-1}P_i}_2$で定義する。
        $P_{i-1}$から$P_i$までの部分曲線から\textcolor{red}{任意に}点$\bm{\xi}_i$をとり、また、$\overrightarrow{P_{i-1}P_i}$の$x_d$軸成分を$(\Delta x_d)_i$とする。
        級数$S(f,\Delta,\{\bm{\xi}_i\}) \coloneqq \sum_{i=1}^n f(\bm{\xi}_i)(\Delta x_d)_i$を考え、これを$f$の$\Delta, \{\bm{\xi}_i\}$に関する\textbf{Riemann和}と呼ぶ。
        分割を任意のやり方で細かくして$|\Delta|\to 0$としたときに$S(f,\Delta,\{\bm{\xi}\})$が収束するとき、これを$\integrate{C}{}{f(\bm{x})}{}{x_d}$と表し、\textbf{$x_d$方向の線積分}という。
    \section{計算法}
        $f$が$C$に関して線積分可能であるとし、パラメータ$t$を用いて$\bm{x}=\bm{\phi}(t)=[\phi_1(t),\dotsc,\phi_n(t)]^\top,\; \bm{\phi} \in C^1(\realNumbers\to\realNumbers^n),\;t:a\to b$と表せる場合は次が成り立つ。
        \[ \integrate{C}{}{f(\bm{x})}{}{x_d} = \integrate{a}{b}{f(\bm{\phi}(t))\phi_d'(t)}{}{t} \]
        その理由を以下で説明する。
        \par
        諸々の定義は「定義」で述べたものを引き続き用いる。
        $t$を$a$から始めて$b$まで増加させると$P_0,P_1,\dotsc$と順に分点に出会う。
        出会った順に対応する$t$を$t_0,t_1,\dotsc$と決めてゆく。
        曲線が自己交差する点では複数の点が重なっていることがあるが、このときは添字の若い順にとる。
        つまり最初の通過で一番若い添字のものが対応し、次の通過で次に若い添字を対応させる。
        平均値の定理より、各$i$に対して適当な$\mu_i\in(t_{i-1},t_i)$が存在して$(\Delta x_d)_i = \phi_d(t_i) - \phi_d(t_{i-1}) = \phi_d'(\mu_i)(t_i-t_{i-1})$を満たす。
        級数$S^{(2)} \coloneqq \sum_{i=1}^n f(\bm{\phi}(\mu_i))\phi_d'(\mu_i)(t_i-t_{i-1})$を考えると、これもまた$f$の$\Delta, \{\bm{\phi}(\mu_i)\}$に関するRiemann和であるから$f$の線積分可能性より、分割を細かくして$|\Delta|\to 0$とするとき$\integrate{C}{}{f(\bm{x})}{}{x_d}$に収束する。
        また、$S^{(2)}$を$t$についての1変数関数$f(\bm{\phi}(t))\phi_d'(t)$のRiemann和と見れば、これは分割を細かくするときに$\integrate{a}{b}{f(\bm{\phi}(t))\phi_d'(t)}{}{t}$に収束する。
        \newline
        \par
        曲線を別のパラメータ表示$\bm{x} = \bm{\psi}(u),\; u:\alpha\to\beta$で表しても、線積分の値は変わらない。
        なぜなら、既に見たように線積分は\textcolor{red}{パラメータ表示に依存せずに定義されている}からである。